\documentclass[a4paper,12pt,titlepage]{article}
% Con este package usa las fuentes Latin Modern Roman (http://www.tug.dk/FontCatalogue/latinmodernroman/)
% En pantalla en el pdf se ven las fuentes mas oscuras (hacer mini-prueba de impresión)
% \usepackage {lmodern}

\usepackage[scaled=1]{couriers}
<<<<<<< HEAD
\usepackage [spanish] {babel} 
=======
% \usepackage [spanish] {babel} 
>>>>>>> ef67754 (Latex: First documentation and examples)
\usepackage [T1]{fontenc}
\usepackage [utf8]{inputenc}
\usepackage {lipsum}
\usepackage [vmargin=1cm,hmargin=2cm,head=26pt,includeheadfoot]{geometry}
\usepackage {fancyhdr}
\usepackage {graphicx}
\usepackage {marginnote}
\usepackage {eso-pic}
\usepackage {booktabs}
\usepackage {float}
\usepackage {tikz}
\usepackage [breaklinks=true,hidelinks]{hyperref}
\usepackage {color}
\usepackage {ifthen}
\usepackage {caption}
\usepackage {parskip}
\usepackage {svn-multi}
\usepackage {verbatim}


\makeatletter
\renewcommand\paragraph{\@startsection{paragraph}{4}{\z@}%
            {-2.5ex\@plus -1ex \@minus -.25ex}%
            {1.25ex \@plus .25ex}%
            {\normalfont\normalsize\bfseries}}
\renewcommand\subparagraph{\@startsection{subparagraph}{5}{\z@}%
            {-2.5ex\@plus -1ex \@minus -.25ex}%
            {1.25ex \@plus .25ex}%
            {\normalfont\normalsize\bfseries}}
\makeatother

\setcounter{secnumdepth}{5}
\setcounter{tocdepth}{5}

%define color in links
\definecolor{colorLinks}{rgb}{0,0,1}
\definecolor{colorGreen}{RGB}{155,187,89}
\definecolor{colorGray}{RGB}{220,220,220}

\hypersetup{
    colorlinks=true,%para colorear los enlaces
    urlcolor=blue,
    linkcolor=blue
}

\fancyhead{}
\fancyfoot{}
\fancyhead[L]{\begin{picture}(0,0) \put(0,0){\includegraphics[scale=.4]{./img/logorealsec.eps}} \end{picture}}
\fancyhead[R]{\begin{picture}(0,0) \put(-90,0){\includegraphics[scale=.4]{./img/encabezado_dcho.eps}} \end{picture}}
\fancyfoot[L]{\marginnote{\begin{picture}(0,0) \put(-525,400){\rotatebox[origin=c]{90}{\begin{tikzpicture}\node [opacity=0.5] (0,0){\tiny La propiedad intelectual de este documento pertenece a \textbf{REALSEC}. Queda prohibida su reproducción, venta, o cesión a terceros.};\end{tikzpicture}}} \end{picture}}[0cm]}
<<<<<<< HEAD
\fancyfoot[C]{Realia Technologies, S.L. \hfill Page \small \thepage \hfill \small \today}
=======
\fancyfoot[C]{Carlos Toro Peñas. \hfill Page \small \thepage \hfill \small \today}
>>>>>>> ef67754 (Latex: First documentation and examples)
\fancyfoot[R]{\small \today}

\pagestyle{fancy}


\begin{document}

%pagina 1: PORTADA
\thispagestyle{empty}
{
	\null\vfill
	\centering\begin{picture}(0,0) \put(-70,170){\includegraphics[scale=.4]{./img/logo_portada.png}} \end{picture}
	\AddToShipoutPicture{\put(90,190){\includegraphics[scale=1]{./img/cover_background.png}}}
<<<<<<< HEAD
	{\Huge \textbf{Universidad Oberta de Catalunya}}\\
	\Large\textbf{\\}
	{\Huge \textbf{Master's degre in Data Science}}\\
	\Large\textbf{\\}
	{\Huge \textbf{Master Thesis}}\\
	{\Huge \textbf{AREA: Specialized fields}}\\
	{
		{
			\Large\textbf{Unsupervised methods applied to Galaxy distribution\\}
		}
	}
	\textcolor{colorGreen}{\\\rule{150mm}{0.3mm}\\}
	\Large\textbf{\\Manual de usuario\\}
	{Rev. 2.0\\}
=======
	{\Large {UNIVERSIDAD OBERTA DE CATALUNYA}}\\
	\Large\textbf{\\}
	{\Large {Master's degree in Data Science}}\\
	\Large\textbf{\\}
	{\Large {MASTER THESIS}}\\
	\Large\textbf{\\}
	{AREA: Specialized fields}\\
	\Large\textbf{\\}
	{
		{
			\Huge\textbf{Unsupervised methods applied to Galaxy distribution\\}
		}
	}
	\textcolor{colorGreen}{\\\rule{150mm}{0.3mm}\\}
	\Large\textbf{\\Author: Carlos Toro Peñas\\}
	\Large\textbf{\\Supervisors: ...\\}
	\Large\textbf{\\Advisors: ...\\}
	\Large\textbf{\\Professor: Laura Ruíz Dern\\}
>>>>>>> ef67754 (Latex: First documentation and examples)
	\vfill
	{\normalsize \textbf{REALIA TECHNOLOGIES, SL} \hfill \small \today }
}

\end{document}